\par\Needspace{7\baselineskip}\bigskip
\begin{center}
\pdfbookmark[2]{Gauss to Olbers, Göttingen, 12 August 1811}{pallas-brief-18110812}
\textsc{Gauss} to \textsc{Olbers}, Göttingen, 12 August 1811.
\end{center}
\noindent\hspace*{\parindent}. . . . .\quad I waited to write to you . . . .\ so that I could add something more about your Pallas. I pass over some things that you will find in one of the next issues of our Gel. Anz.\HistoricalNoteEight{Vol. VI, p. 327.]} (I believe the one dated 17 Aug.), so as to mention only some of the general perturbation results. In the first calculation I finished working out the perturbations of the inclination and of the longitude of the line of nodes (the least laborious part), and I also checked the former at the 7 oppositions (for which my earlier work gives the elements). Poor agreement at first, until I found that 2 equations had the wrong sign. They now agree to my satisfaction. I shall also finish checking the \(\AscendingNode\) in a few days, and then I intend to tackle the eccentricity, about a month's work. To give you some idea of the results, I enclose the 40 perturbation equations for the inclination. I have retained only what exceeded one arcsecond. A few small correction terms depending on 9 times and 10 times Jupiter's longitude might still be added\HistoricalNoteDouble{The terms of \(\delta i\) given on p. 518 of this volume follow here.]}.\quad . . . . .
\par
The 80 equations for inclination and \(\AscendingNode\) could be combined into 40 for latitude, but I do not think anything is gained, for if one allows the elements themselves to be perturbed, one can without hesitation treat the same perturbed elements as valid for several months. One therefore needs to calculate the perturbations for 6 elements only once a year (perhaps about 300—400 equations in all). If, however, one applies only secular changes to the elements and applies the periodic changes to latitude, longitude and radius vector (perhaps close to 200 in all), these would presumably have to be calculated for at least 6 different positions in a year to allow interpolation. But in the \SourceLeaf{427}future one can proceed as one wishes, provided one first has all the perturbations \textls[180]{in some form}.
\par
